\subsection{全纯函数的芽}

设 E 与 F 是复 Banach 空间。我们回顾（cf. VAR, R1, p.~26, 3.2.1, p.~22, 3.1 及 p.~88, App.）：定义在 E 的开子集 U 上、取值于 F 中的全纯映射是这样的映射 \(f \colon \mathrm{U} \to \mathrm{F}\)：对一切 \(x \in \mathrm{U}\)，存在一个收敛级数

\[
f _ {x} = \sum_ {k \geqslant 0} f _ {x, k}
\]

它满足 \(f(x + y) = f_x(y)\)（对充分接近 0 的一切 \(y \in E\)），其中 \(f_{x,k} \colon E \to C\) 是 E 上的、取值于 F 中的 k 次连续齐次多项式，即形如

\[
f _ {x, k} (y) = \widetilde {f} _ {x, k} (y, \ldots , y)
\]

其中 \(\widetilde{f}_{x,k}\colon\mathrm{E}^{k}\to\mathrm{F}\) 是一个 \(k\) 线性连续映射。我们用 \(\mathcal{O}(\mathrm{U};\mathrm{F})\) 表示 U 上取值于 F 中的全纯函数所成的复向量空间，赋以紧收敛拓扑。它是一个局部凸拓扑向量空间，其拓扑由半范数 \(f\mapsto\sup_{z\in\mathrm{K}}\|f(z)\|\) 定义，这里 K 取遍 U 的紧子集所成的集合。

设 G 是一个复 Banach 空间，V 是 G 的一个开子集。对每个全纯映射 \(\varphi\colon\mathrm{V}\to\mathrm{U}\)，映射 \(\varphi^{*}\colon f\mapsto f\circ\varphi\) 是从 \(\mathcal{O}(\mathrm{U};\mathrm{F})\) 到 \(\mathcal{O}(\mathrm{V};\mathrm{F})\) 中的连续线性映射。

若 H 是一个复 Banach 空间，\(\varphi\colon\mathrm{F}\to\mathrm{H}\) 是一个连续线性映射，则映射 \(f\mapsto\varphi\circ f\) 是从 \(\mathcal{O}(\mathrm{U};\mathrm{F})\) 到 \(\mathcal{O}(\mathrm{U};\mathrm{H})\) 中的连续线性映射，记作 \(\varphi_{*}\)。

设 n 是一个自然数，并令 \(E = C^{n}\)。设 K 是 \(C^{n}\) 的一个紧子集，U 是 K 的开邻域组成的递降过滤集。若 U、\(U' \in U\) 且 \(U' \subset U\)，则从 \(\mathcal{O}(U; F)\) 到 \(\mathcal{O}(U'; F)\) 中的函数限制映射是连续的。诸空间 \(\mathcal{O}(\mathrm{U};\mathrm{F})\) 关于这些映射的归纳极限记作 \(\mathcal{O}(\mathrm{K};\mathrm{F})\)。\(\mathcal{O}(\mathrm{K};\mathrm{F})\) 的元素称为 K 的邻域中的、取值于 F 的全纯函数芽。

空间 \(\mathcal{O}(\mathrm{K};\mathrm{F})\) 赋以诸 \(\mathcal{O}(\mathrm{U};\mathrm{F})\) 的局部凸拓扑的归纳极限拓扑。设 X 是一个局部凸拓扑向量空间，\(\varphi\colon\mathcal{O}(\mathrm{K};\mathrm{F})\to\mathrm{X}\) 是一个映射。对 K 的每个开邻域 U，映射 \(\mathcal{O}(\mathrm{U};\mathrm{F})\to\mathrm{X}\)（由 \(\varphi\) 与典范映射 \(\mathcal{O}(\mathrm{U};\mathrm{F})\to\mathcal{O}(\mathrm{K};\mathrm{F})\) 复合而得到）记作 \(\varphi^{\mathrm{U}}\)。映射 \(\varphi\) 连续当且仅当对每个 U，\(\varphi^{\mathrm{U}}\) 连续 (EVT, II, p.~29, prop. 5, (iii))。

设 m 是一个自然数。设 L 是 \(C^{m}\) 的一个紧子集，V 是 L 的一个开邻域。设 \(\varphi\colon V \to C^{n}\) 是一个满足 \(\varphi(L) \subset K\) 的全纯映射。连续线性映射

\[
\mathcal {O} (\mathrm{U}; \mathrm{F}) \xrightarrow {\varphi^ {*}} \mathcal {O} (\stackrel {- 1} {\varphi} (\mathrm{U}); \mathrm{F}) \xrightarrow {\varphi^ {- 1} \stackrel {- 1} {\varphi} (\mathrm{U})} \mathcal {O} (\mathrm{L}; \mathrm{F})
\]

对 K 的每个开邻域 U，诱导出一个连续线性映射 \(\varphi^{*}\colon\mathcal{O}(\mathrm{K};\mathrm{F})\to\mathcal{O}(\mathrm{L};\mathrm{F})\) (loc. cit.)。

设 H 是一个复 Banach 空间，\(\varphi\colon\mathrm{F}\to\mathrm{H}\) 是一个连续线性映射。连续线性映射

\[
\mathcal {O} (\mathrm{U}; \mathrm{F}) \xrightarrow {\varphi_ {*}} \mathcal {O} (\mathrm{U}; \mathrm{H}) \xrightarrow {\varphi^ {\mathrm{U}}} \mathcal {O} (\mathrm{K}; \mathrm{F})
\]

（其中 U 取遍 K 在 \(C^{n}\) 中的开邻域所成的集合）诱导出一个连续线性映射 \(\varphi_{*}\)（从 \(\mathcal{O}(K;F)\) 到 \(\mathcal{O}(K;H)\) 中）(loc. cit.)。我们有时将记 \(\varphi \circ f = \varphi_{*}(f)\)。

对 K 的每个开邻域 U，在 K 上的限制是一个连续线性映射 \(\mathcal{O}(\mathrm{U};\mathrm{F})\to \mathcal{C}(\mathrm{K};\mathrm{F})\)；这些映射诱导出一个连续线性映射 \(\mathcal{O}(\mathrm{K};\mathrm{F})\to \mathcal{C}(\mathrm{K};\mathrm{F})\)，称为 K 上全纯函数芽的赋值。

设 A 是一个复幺 Banach 代数。空间 \(\mathcal{O}(\mathrm{U};\mathrm{A})\) 与 \(\mathcal{O}(\mathrm{K};\mathrm{A})\) 是幺代数。若 \(A \neq \{0\}\)，则可以把 \(\mathcal{O}(\mathrm{U};\mathbf{C})\)（相应地，\(\mathcal{O}(\mathrm{K};\mathbf{C})\)）典范地视为子代数 \(\mathcal{O}(\mathrm{U};\mathbf{C}) \cdot 1\)（\(\mathcal{O}(\mathrm{U};\mathrm{A})\) 的子代数）（相应地，子代数 \(\mathcal{O}(\mathrm{K};\mathbf{C}) \cdot 1\)（\(\mathcal{O}(\mathrm{K};\mathrm{A})\) 的子代数））。我们将令 \(\mathcal{O}(\mathrm{U}) = \mathcal{O}(\mathrm{U};\mathbf{C})\) 及 \(\mathcal{O}(\mathrm{K}) = \mathcal{O}(\mathrm{K};\mathbf{C})\)。

